% Seminar 4 - Page 5. Fragment for
% % Seminar 4 - Page 5. Fragment for
% % Seminar 4 - Page 5. Fragment for
% % Seminar 4 - Page 5. Fragment for
% \input{seminars/pages/seminar4-pagina5-en}; no preamble or forced page break.
% EDITORIAL NOTES:
% - The opening is a counterexample to a Cantor-Bernstein-type assertion
%   for groups: each group can be isomorphic to a subgroup of the other,
%   without the groups being isomorphic. The handwritten implication
%   ending in G_1 not isomorphic to G_2 is interpreted in this context,
%   not as a universal implication.
% - The tuples are handwritten as (q_0,q_1,...,q_n), while the displayed
%   shift phi(q_0,...,q_n)=(0,q_0,...) only gives the intended group
%   isomorphism when the coordinates continue indefinitely. This
%   indexing ambiguity is NOT silently repaired in the displayed text.
%   For finite fixed n, Q^(n+1) is not isomorphic to 0 x Q^n.
% - The reason G_1 and G_2 are not isomorphic is not proved in the notes.
%   For the infinite-sequence reading, G_1 is divisible, whereas G_2
%   (with its integral first coordinate) is not.
% - The final line lists orders 4,5,6,7, but spells out only the groups
%   of orders 4,5,6. No classification for order 7 is written down.
% - An ink correction precedes the Z_6/S_3 line; only the legible
%   group names are transcribed. Very faint show-through is omitted.

\subsection*{Seminar 4 - Page 5}

\noindent Let $G_1,G_2$ be groups. The manuscript considers the situation
\[
\begin{gathered}
 G_1\text{ is isomorphic to a subgroup of }G_2,\\
 G_2\text{ is isomorphic to a subgroup of }G_1,
\end{gathered}
\qquad\text{but}\qquad G_1\not\cong G_2.
\]

\medskip
\noindent Define
\[
 G_1:=\bigl\{x=(q_0,q_1,\ldots,q_n)\mid q_i\in\mathbb Q\bigr\}.
\]
Then $(G_1,+)$ is a group. For
$y=(u_0,u_1,\ldots,u_n)$,
\[
 x+y=(q_0+u_0,\ldots,q_n+u_n).
\]

\noindent Consider also
\[
 \begin{aligned}
 G_2:=\bigl\{\,&x=(q_0,q_1,\ldots,q_n)\mid
     q_0\in\mathbb Z,\quad q_i\in\mathbb Q\ 
     (1\leq i\leq n)\,\bigr\}\subseteq G_1.
 \end{aligned}
\]
The set $(G_2,+)$ is a group and $G_2\leq G_1$
($G_2$ is a subgroup of $G_1$).

\medskip
\noindent Put
\[
 H:=\bigl\{x=(0,q_1,\ldots,q_n)\mid
     q_i\in\mathbb Q\ (1\leq i\leq n)\bigr\}.
\]
Then $H\leq G_2$. The notes give the map
\[
 \begin{aligned}
 \varphi:G_1&\longrightarrow H,\\
 \varphi\bigl((q_0,\ldots,q_n)\bigr)&=(0,q_0,\ldots),
 \end{aligned}
\]
and state that $\varphi$ is a group isomorphism.

\medskip
\noindent\textbf{Conclusion (as in the manuscript).}
$G_1$ and $G_2$ are not isomorphic groups.

\medskip
\noindent\underline{Groups of small order.}
The manuscript mentions $|G|\in\{4,5,6,7\}$ and records
\[
\begin{array}{r@{\quad\Longrightarrow\quad}l}
 |G|=4 & G\cong\mathbb Z_4\quad\text{or}\quad
               \mathbb Z_2\times\mathbb Z_2,\\[2pt]
 |G|=5 & G\cong\mathbb Z_5,\\[2pt]
 |G|=6 & G\cong\mathbb Z_6\quad\text{or}\quad S_3.
\end{array}
\]
; no preamble or forced page break.
% EDITORIAL NOTES:
% - The opening is a counterexample to a Cantor-Bernstein-type assertion
%   for groups: each group can be isomorphic to a subgroup of the other,
%   without the groups being isomorphic. The handwritten implication
%   ending in G_1 not isomorphic to G_2 is interpreted in this context,
%   not as a universal implication.
% - The tuples are handwritten as (q_0,q_1,...,q_n), while the displayed
%   shift phi(q_0,...,q_n)=(0,q_0,...) only gives the intended group
%   isomorphism when the coordinates continue indefinitely. This
%   indexing ambiguity is NOT silently repaired in the displayed text.
%   For finite fixed n, Q^(n+1) is not isomorphic to 0 x Q^n.
% - The reason G_1 and G_2 are not isomorphic is not proved in the notes.
%   For the infinite-sequence reading, G_1 is divisible, whereas G_2
%   (with its integral first coordinate) is not.
% - The final line lists orders 4,5,6,7, but spells out only the groups
%   of orders 4,5,6. No classification for order 7 is written down.
% - An ink correction precedes the Z_6/S_3 line; only the legible
%   group names are transcribed. Very faint show-through is omitted.

\subsection*{Seminar 4 - Page 5}

\noindent Let $G_1,G_2$ be groups. The manuscript considers the situation
\[
\begin{gathered}
 G_1\text{ is isomorphic to a subgroup of }G_2,\\
 G_2\text{ is isomorphic to a subgroup of }G_1,
\end{gathered}
\qquad\text{but}\qquad G_1\not\cong G_2.
\]

\medskip
\noindent Define
\[
 G_1:=\bigl\{x=(q_0,q_1,\ldots,q_n)\mid q_i\in\mathbb Q\bigr\}.
\]
Then $(G_1,+)$ is a group. For
$y=(u_0,u_1,\ldots,u_n)$,
\[
 x+y=(q_0+u_0,\ldots,q_n+u_n).
\]

\noindent Consider also
\[
 \begin{aligned}
 G_2:=\bigl\{\,&x=(q_0,q_1,\ldots,q_n)\mid
     q_0\in\mathbb Z,\quad q_i\in\mathbb Q\ 
     (1\leq i\leq n)\,\bigr\}\subseteq G_1.
 \end{aligned}
\]
The set $(G_2,+)$ is a group and $G_2\leq G_1$
($G_2$ is a subgroup of $G_1$).

\medskip
\noindent Put
\[
 H:=\bigl\{x=(0,q_1,\ldots,q_n)\mid
     q_i\in\mathbb Q\ (1\leq i\leq n)\bigr\}.
\]
Then $H\leq G_2$. The notes give the map
\[
 \begin{aligned}
 \varphi:G_1&\longrightarrow H,\\
 \varphi\bigl((q_0,\ldots,q_n)\bigr)&=(0,q_0,\ldots),
 \end{aligned}
\]
and state that $\varphi$ is a group isomorphism.

\medskip
\noindent\textbf{Conclusion (as in the manuscript).}
$G_1$ and $G_2$ are not isomorphic groups.

\medskip
\noindent\underline{Groups of small order.}
The manuscript mentions $|G|\in\{4,5,6,7\}$ and records
\[
\begin{array}{r@{\quad\Longrightarrow\quad}l}
 |G|=4 & G\cong\mathbb Z_4\quad\text{or}\quad
               \mathbb Z_2\times\mathbb Z_2,\\[2pt]
 |G|=5 & G\cong\mathbb Z_5,\\[2pt]
 |G|=6 & G\cong\mathbb Z_6\quad\text{or}\quad S_3.
\end{array}
\]
; no preamble or forced page break.
% EDITORIAL NOTES:
% - The opening is a counterexample to a Cantor-Bernstein-type assertion
%   for groups: each group can be isomorphic to a subgroup of the other,
%   without the groups being isomorphic. The handwritten implication
%   ending in G_1 not isomorphic to G_2 is interpreted in this context,
%   not as a universal implication.
% - The tuples are handwritten as (q_0,q_1,...,q_n), while the displayed
%   shift phi(q_0,...,q_n)=(0,q_0,...) only gives the intended group
%   isomorphism when the coordinates continue indefinitely. This
%   indexing ambiguity is NOT silently repaired in the displayed text.
%   For finite fixed n, Q^(n+1) is not isomorphic to 0 x Q^n.
% - The reason G_1 and G_2 are not isomorphic is not proved in the notes.
%   For the infinite-sequence reading, G_1 is divisible, whereas G_2
%   (with its integral first coordinate) is not.
% - The final line lists orders 4,5,6,7, but spells out only the groups
%   of orders 4,5,6. No classification for order 7 is written down.
% - An ink correction precedes the Z_6/S_3 line; only the legible
%   group names are transcribed. Very faint show-through is omitted.

\subsection*{Seminar 4 - Page 5}

\noindent Let $G_1,G_2$ be groups. The manuscript considers the situation
\[
\begin{gathered}
 G_1\text{ is isomorphic to a subgroup of }G_2,\\
 G_2\text{ is isomorphic to a subgroup of }G_1,
\end{gathered}
\qquad\text{but}\qquad G_1\not\cong G_2.
\]

\medskip
\noindent Define
\[
 G_1:=\bigl\{x=(q_0,q_1,\ldots,q_n)\mid q_i\in\mathbb Q\bigr\}.
\]
Then $(G_1,+)$ is a group. For
$y=(u_0,u_1,\ldots,u_n)$,
\[
 x+y=(q_0+u_0,\ldots,q_n+u_n).
\]

\noindent Consider also
\[
 \begin{aligned}
 G_2:=\bigl\{\,&x=(q_0,q_1,\ldots,q_n)\mid
     q_0\in\mathbb Z,\quad q_i\in\mathbb Q\ 
     (1\leq i\leq n)\,\bigr\}\subseteq G_1.
 \end{aligned}
\]
The set $(G_2,+)$ is a group and $G_2\leq G_1$
($G_2$ is a subgroup of $G_1$).

\medskip
\noindent Put
\[
 H:=\bigl\{x=(0,q_1,\ldots,q_n)\mid
     q_i\in\mathbb Q\ (1\leq i\leq n)\bigr\}.
\]
Then $H\leq G_2$. The notes give the map
\[
 \begin{aligned}
 \varphi:G_1&\longrightarrow H,\\
 \varphi\bigl((q_0,\ldots,q_n)\bigr)&=(0,q_0,\ldots),
 \end{aligned}
\]
and state that $\varphi$ is a group isomorphism.

\medskip
\noindent\textbf{Conclusion (as in the manuscript).}
$G_1$ and $G_2$ are not isomorphic groups.

\medskip
\noindent\underline{Groups of small order.}
The manuscript mentions $|G|\in\{4,5,6,7\}$ and records
\[
\begin{array}{r@{\quad\Longrightarrow\quad}l}
 |G|=4 & G\cong\mathbb Z_4\quad\text{or}\quad
               \mathbb Z_2\times\mathbb Z_2,\\[2pt]
 |G|=5 & G\cong\mathbb Z_5,\\[2pt]
 |G|=6 & G\cong\mathbb Z_6\quad\text{or}\quad S_3.
\end{array}
\]
; no preamble or forced page break.
% EDITORIAL NOTES:
% - The opening is a counterexample to a Cantor-Bernstein-type assertion
%   for groups: each group can be isomorphic to a subgroup of the other,
%   without the groups being isomorphic. The handwritten implication
%   ending in G_1 not isomorphic to G_2 is interpreted in this context,
%   not as a universal implication.
% - The tuples are handwritten as (q_0,q_1,...,q_n), while the displayed
%   shift phi(q_0,...,q_n)=(0,q_0,...) only gives the intended group
%   isomorphism when the coordinates continue indefinitely. This
%   indexing ambiguity is NOT silently repaired in the displayed text.
%   For finite fixed n, Q^(n+1) is not isomorphic to 0 x Q^n.
% - The reason G_1 and G_2 are not isomorphic is not proved in the notes.
%   For the infinite-sequence reading, G_1 is divisible, whereas G_2
%   (with its integral first coordinate) is not.
% - The final line lists orders 4,5,6,7, but spells out only the groups
%   of orders 4,5,6. No classification for order 7 is written down.
% - An ink correction precedes the Z_6/S_3 line; only the legible
%   group names are transcribed. Very faint show-through is omitted.

\subsection*{Seminar 4 - Page 5}

\noindent Let $G_1,G_2$ be groups. The manuscript considers the situation
\[
\begin{gathered}
 G_1\text{ is isomorphic to a subgroup of }G_2,\\
 G_2\text{ is isomorphic to a subgroup of }G_1,
\end{gathered}
\qquad\text{but}\qquad G_1\not\cong G_2.
\]

\medskip
\noindent Define
\[
 G_1:=\bigl\{x=(q_0,q_1,\ldots,q_n)\mid q_i\in\mathbb Q\bigr\}.
\]
Then $(G_1,+)$ is a group. For
$y=(u_0,u_1,\ldots,u_n)$,
\[
 x+y=(q_0+u_0,\ldots,q_n+u_n).
\]

\noindent Consider also
\[
 \begin{aligned}
 G_2:=\bigl\{\,&x=(q_0,q_1,\ldots,q_n)\mid
     q_0\in\mathbb Z,\quad q_i\in\mathbb Q\ 
     (1\leq i\leq n)\,\bigr\}\subseteq G_1.
 \end{aligned}
\]
The set $(G_2,+)$ is a group and $G_2\leq G_1$
($G_2$ is a subgroup of $G_1$).

\medskip
\noindent Put
\[
 H:=\bigl\{x=(0,q_1,\ldots,q_n)\mid
     q_i\in\mathbb Q\ (1\leq i\leq n)\bigr\}.
\]
Then $H\leq G_2$. The notes give the map
\[
 \begin{aligned}
 \varphi:G_1&\longrightarrow H,\\
 \varphi\bigl((q_0,\ldots,q_n)\bigr)&=(0,q_0,\ldots),
 \end{aligned}
\]
and state that $\varphi$ is a group isomorphism.

\medskip
\noindent\textbf{Conclusion (as in the manuscript).}
$G_1$ and $G_2$ are not isomorphic groups.

\medskip
\noindent\underline{Groups of small order.}
The manuscript mentions $|G|\in\{4,5,6,7\}$ and records
\[
\begin{array}{r@{\quad\Longrightarrow\quad}l}
 |G|=4 & G\cong\mathbb Z_4\quad\text{or}\quad
               \mathbb Z_2\times\mathbb Z_2,\\[2pt]
 |G|=5 & G\cong\mathbb Z_5,\\[2pt]
 |G|=6 & G\cong\mathbb Z_6\quad\text{or}\quad S_3.
\end{array}
\]
